\documentclass[10pt,a4paper]{amsart}
\usepackage[margin=19mm]{geometry}
\usepackage{amsmath,amssymb,mathtools}
\usepackage[scr=rsfs,cal=euler]{mathalfa}
\usepackage{xfrac}
\usepackage[unicode,hidelinks,bookmarks=false]{hyperref}
\newcommand{\ie}{i.e.\ }
\newcommand{\cf}{cf.\ }
\newcommand{\resp}{respectively\ }
\newcommand{\Subsection}{\subsection}
\providecommand{\nobreakdash}{\nobreak\hbox{-}}
\setlength{\emergencystretch}{2em}
\multlinegap0pt
\usepackage{amssymb}
\usepackage{enumerate}
\usepackage{mathtools}
\usepackage{xcolor}
\newtheorem{theorem}{Theorem}
\newcommand\NN{\mathbb{N}}
\newcommand\QQ{\mathbb{Q}}
\renewcommand{\leq}{\leqslant}
\title{Ratio of sum of digits functions in two bases}
\author{Pascal Jelinek}
\date{Source: arXiv:2608.14241v2 (CC BY 4.0)}
\begin{document}
\maketitle

\noindent\emph{Source and licence.} Ratio of sum of digits functions in two bases. Version arXiv:2608.14241v2 (CC BY 4.0), distributed by arXiv under the Creative Commons Attribution 4.0 International licence, http://creativecommons.org/licenses/by/4.0/, as recorded in the arXiv licence metadata for this exact version and shown on the abstract page https://arxiv.org/abs/2608.14241v2. Full licence notice: \texttt{licenses/arxiv-2608.14241-CC-BY-4.0.txt}.\par\smallskip

\noindent\textit{Editorial scope.} Counted unit: the article's theorem main\_thm\_dep (source0.tex lines 206--220). The article's complete original proof of it is delivered with it (source0.tex lines 1472--1547, one \texttt{proof} environment of 76 lines, which the article places away from the statement). No mathematics is written, repaired or re-derived: the statement and the proof are the article's own lines, in the article's own notation, and no retained line is substituted or re-flowed.  No further source-local context is needed: every cross-reference of every retained line resolves inside the two delivered spans.  Nothing was omitted from any retained span, so the package carries no omission record.  No retained line contains a citation, so the wrapper carries no bibliography and no bibliography entry.  No retained line contains a figure environment, an image-inclusion command or drawing code, so no image is delivered and the package declares no asset dependency.  Editorial formatting only, in addition: an AMS \texttt{amsart} preamble that restates the article's own declarations and macros used by the retained lines, namely the environments \texttt{theorem} on separate counters, together with the standard packages those lines need.  The retained environments print in this document as Theorem 1 in order of appearance, whereas the article prints the same items with the numbering of its own full document.  The complete native source of the article as distributed in the arXiv e-print for this exact version is bundled unchanged as \texttt{sources/207651/source0.tex}.
\begin{theorem}\label{main_thm_dep}
	Let $1<q_1,q_2$ be two distinct multiplicatively dependent numbers, say $q_1=b^k$ and $q_2=b^{\ell}$ for some coprime natural numbers $k,\ell$  and let $r\in \QQ^+$. Then we have that there are infinitely many $n\in \NN$ such that
	\[ 
		\frac{s_{q_1}(n)}{s_{q_2}(n)}=r
	\]
	as long as
	\[ 
		\frac{1}{b^{\ell-1}}\leq r \leq b^{k-1}.
	\]
More precisely, we have that
	\[
		\#\left\{1<n\leq N: \frac{s_{q_1}(n)}{s_{q_2}(n)}=r\right\} \gg N^{c},
	\]
	where $0<c\leq 1$ is an effectively computable constant, which depends on $k$ and $\ell$ 
\end{theorem}

\begin{proof}[Proof of Theorem~\ref{main_thm_dep}]
	The key to the proof is the following observation:\\
	Since $q_1=b^k$ and $q_2=b^{\ell}$ we see that if we group the digits from the right in blocks of $k$ digits in base $b$, each such block corresponds to a digit in base $q_1$. Similarly, if we group the digits from the right in blocks of $\ell$ digits in base $b$, each such block corresponds to a digit in base $q_2$.\\
	Our first aim is to show that $\frac{1}{b^{\ell-1}}$ and $b^{k-1}$ can each be attained infinitely many times. For this observe the following identities for the digit sums of $b^x$ for some $x\in \NN$, and $a < b$ a non-negative integer:
	\[
	s_{q_1}(ab^x)=ab^{x\bmod k} \text{ and } s_{q_2}(ab^x)=ab^{x\bmod \ell}
	\]
	We can now use the Chinese remainder theorem by our assumption that $k$ and $\ell$ are coprime, to find infinitely many $x$, such that $x\equiv k-1 \bmod k$ and $x~\equiv 0 \bmod \ell$, which gives for any $a< b$ that
	\[
	\frac{s_{q_1}(ab^x)}{s_{q_2}(ab^x)}=\frac{ab^{k-1}}{ab^0}=b^{k-1}.
	\]
	Similarly, we can find infinitely many $x$ such that $x \equiv 0 \bmod k$ and $x\equiv \ell -1 \bmod \ell$ which corresponds to the ratio
	\[
	\frac{s_{q_1}(ab^x)}{s_{q_2}(ab^x)}=\frac{1}{b^{\ell-1}}.
	\]
	To see that these values are indeed the maximum and the minimum, respectively, we write $n=\sum_{x\in \NN}a_xb^x$ in the base $b$ expansion. By the fact that $k$ digits in base $b$ correspond to a single digit in base $q_1$ we have that 
	\[
	s_{q_1}(n) = s_{q_1}\left(\sum_{x\in \NN}a_xb^x\right) = \sum_{x\in \NN}a_xs_{q_1}(b^x)=\sum_{x\in \NN}a_xb^{x\bmod k}.
	\]
	Similarly, in base $q_2$ we have that
	\[
	s_{q_2}(n) =\sum_{x\in \NN}a_xb^{x\bmod \ell}.
	\]
	Therefore, we have that
	\begin{align*}
		\frac{s_{q_1}(n)}{s_{q_2}(n)} &= \frac{\sum_{x\in \NN}a_xb^{x\bmod k}}{\sum_{x\in \NN}a_xb^{x\bmod \ell}}.
	\end{align*}
	One can easily show by induction that
	\begin{align*}
		\min_{\substack{x\in\NN\\a_x\neq 0}}\left(\frac{b^{x\bmod k}}{b^{x\bmod \ell}}\right) \leq \frac{\sum_{x\in \NN}a_xb^{x\bmod k}}{\sum_{x\in \NN}a_xb^{x\bmod \ell}} \leq \max_{\substack{x\in\NN\\a_x\neq 0}}\left(\frac{b^{x\bmod k}}{b^{x\bmod \ell}}\right).
	\end{align*}
	Therefore, we see that 
	\[ 
	\frac{1}{b^{\ell-1}}\leq r \leq b^{k-1}.
	\]
	It only remains to prove that for any such rational number $r=u/v$, with $u,v$ coprime, we can find an $n\in \NN$ such that
	\[
	\frac{s_{q_1}(n)}{s_{q_2}(n)}=\frac uv
	\]
	Let $x>0$ be minimal such that $x\equiv k-1 \bmod k$ and $x\equiv 0 \bmod \ell$, and let $x_i=x+ik\ell$. Similarly, let $y>0$ be minimal such that $y\equiv 0 \bmod k$ and $y\equiv \ell-1 \bmod \ell$, and let $y_i=y+ik\ell$. Consider
	\[
	n=\sum_{i=1}^{s}b^{x_i}+\sum_{i=s+1}^{s+t}b^{y_i}
	\]
	for some $s,t>0$. Our observations above give that $s_{q_1}(n)= sb^{k-1}+tb^{0}$ and $s_{q_2}(n)= sb^{0}+tb^{\ell-1}$. Therefore, we are trying to solve the following equation for $s,t>0$:
	\[
	\frac{sb^{k-1}+t}{s+tb^{\ell-1}}=\frac uv
	\]
	Rearranging gives that
	\[
	s(vb^{k-1}-u)=t(ub^{\ell-1}-v)\,.
	\]
	We see that the brackets are both positive if (and only if)
	\[ 
	\frac{1}{b^{\ell-1}}\leq \frac uv \leq b^{k-1}\,.
	\]
	Therefore, for any such $u,v$ the equation has infinitely many solutions for $s,t>0$. Hence, recalling that $n=\sum_{i=1}^{s}b^{x_i}+\sum_{i=s+1}^{s+t}b^{y_i}$ we see that
	\[
	\frac{s_{q_1}(n)}{s_{q_2}(n)}=\frac uv.
	\]
	To get a quantitative lower bound, we consider 
	\[
	n=a+\sum_{i=r}^{r+s}b^{x_i}+\sum_{i=r+s+1}^{s+r+t}b^{y_i}
	\]
	where $r$ is minimal such that $a<b^{x_i}$. We say that $s_{q_1}(a)=e$ and $s_{q_2}(a)=f$ and we again want to solve for $s,t>0$ such that
	\[
	\frac{sb^{k-1}+t+e}{s+tb^{\ell-1}+f}=\frac uv.
	\]
	Rearranging the equation above, we get that
	\[
	s(vb^{k-1}-u)-t(ub^{\ell-1}-v)=f(ub^{\ell-1}-v)-e(vb^{k-1}-u)\,.
	\]
	If we assume that $a\leq b^M$, we have that $e,f\leq cM$, for some constant $c$. Therefore, we can find a solution for $s,t$ to the above equation if $s,t$ are of the size $\sim cM$ and hence $n\leq b^{c_1M}$. 
	
	In particular, this shows that given all integers up to $N$, we have at least $N^c$ many numbers $n$ such that
	\[\frac{s_{q_1}(n)}{s_{q_2}(n)}=\frac uv.\]
\end{proof}

\end{document}
